\section{Orientações do trabalho}

Esta seção resume, em formato de checklist, como cada uma das etapas solicitadas nas orientações do trabalho foi contemplada neste relatório. As atividades foram adaptadas ao conjunto de dados \texttt{iris}, utilizado aqui como exemplo.

\begin{longtable}[]{@{}p{0.5\linewidth}p{0.42\linewidth}@{}}
\caption{Checklist de aderência às orientações do trabalho.}\tabularnewline
\toprule\noalign{}
Item & Seção do relatório \\
\midrule\noalign{}
\endfirsthead
\toprule\noalign{}
Item & Seção do relatório \\
\midrule\noalign{}
\endhead
\bottomrule\noalign{}
\endlastfoot
1. Problema, características dos dados e objetivo da análise & Introdução \\
2. Descrição da base (n, variáveis, significado, unidade) & Descrição da base de dados (Tabela 1) \\
3. Análise exploratória (dados ausentes, outliers, objetivos) & Análise exploratória \\
4. Análises uni/bivariadas: correlação, multicolinearidade, box-plot & Análise de correlação; Multicolinearidade: o VIF \\
5. Modelo inicial com todas as variáveis explicativas & Multicolinearidade: o VIF (Tabela 3) \\
6. Seleção do melhor modelo por critério estatístico & Seleção do modelo final (AIC, BIC, teste F, stepwise) \\
7. Verificação e resolução de violações de pressupostos & Diagnóstico do modelo selecionado; Impacto da remoção \\
8a. IC para o valor médio esperado da resposta, com interpretação & Previsões: Intervalo de confiança \\
8b. Previsão pontual e intervalo de predição, com interpretação & Previsões: Intervalo de predição \\
\end{longtable}
